% ===== §18 Worked Example: Cultivating the Perception of Contingency =====
\section{A Worked Example: The Cultivation of Contingency}
\label{sec:worked-example}

\noindent
The four movements have been given their theory, and the theory will remain an abstraction until it is shown doing the one thing it was built to do, which is to be practised. This section takes a single capacity, the aesthetic perception of contingency, and walks it through all four steps, naming for each the practice that cultivates it, the failure that shadows it, and the condition under which it becomes the principal site. The capacity is chosen deliberately. A companion to this series argued that the perception of contingency, the noticing of what need not have happened, is the aesthetic core of the field of travel, the openness in which the relational subject is generated; but that argument said what the perception is and not how it is trained, and the omission is the one this section repairs. What follows is therefore not an illustration laid beside the theory but the theory brought to the scale at which a person could begin tomorrow morning, and it carries one constraint that the political economy section will defend: every practice named here must be a re-meeting of an act the day already contains, and not an added exercise requiring leisure the worn-down do not have, since an aesthetic education affordable only to the unhurried would betray its own thesis.

\subsection{Perceiving contingency}
\label{sub:we-perceiving}

The first step is to see the contingent at all, to register the small unplanned thing that the day would otherwise absorb without trace: the unexpected turn of a sentence, the way a light fell that will not fall so again, a reaction from the other that the script did not predict. The practice is the insertion of a half-second before a habitual act, a refusal to begin the planned motion until the field as it actually is this morning has been let to appear, the other's tone today, the particular disorder of the table, the thing that is not as it was yesterday. This is the emptying-to-await of the Daoist voice made into a repeatable micro-act, a gap held open between intention and the field so that the prior does not flatten the morning into a copy of every prior morning. The failure that shadows it is the flattening itself, the strong expectation that explains the field in advance, the silent certainty that she is as she always is at this hour, under which the contingent is registered as nothing because it was predicted as everything. The sign of whether the step is alive is simple and merciless: whether one can still be touched, inside the most familiar hour, by something small that need not have been.

This is the principal site when a relation has grown dull. Dullness is not the absence of activity, since the dull relation still acts, still greets and sets the table and shares by habit; it is the cessation of seeing, the field grown so predicted that no contingency gets through, and the labour that matters then is not more doing but the recovery of sight, the reopening of the half-second gap that habit has closed. To work on sharing or reproduction while perception has died is to elaborate a cycle that no longer has an intake, and the dialectical discipline of the model is what tells the practitioner, here, that the gap is where the work belongs.

\subsection{Receiving contingency}
\label{sub:we-receiving}

To perceive the contingent is not yet to let it in. The second step is to allow the unexpected thing to alter the morning rather than to correct it back into the plan, and the practice is the suppression of the reflex by which the unforeseen is treated as a deviation to be managed, the answer that did not fit folded back toward the answer that would have, the unplanned mood received as an obstacle to the day's intentions rather than as a thing that may rewrite them. This is the letting-be of releasement and the mirror that does not see things off, made into the discipline of not putting the contingent back in its place, which the reception section showed to be, at its depth, the refusal to reduce the other to the same. The failure that shadows it is defensive gating, the meeting of the unexpected as interference, its swift return to order before it can change anything, the morning's plan defended against the morning. The sign of the step's health is whether what was not foreseen is permitted to revise the valuation it entered, whether the day can be, in some small part, rewritten by what it did not expect.

This is the principal site when a relation has lately been wounded, because injury is exactly what raises the gate. After hurt the defended posture is highest, the unexpected is met as threat by default, and the contingent that might have refreshed the field cannot get past the guard; no quantity of perceiving, no eagerness to share, substitutes for the lowering of that gate, which is why the reception section insisted that the wounded relation's work is reception and not the steps around it. The mercy the political economy section will urge belongs here too: the gate that injury has shut does not open on command, and the practice is the patient, undemanded lowering of a guard, not the forcing of a welcome the wound has made impossible.

\subsection{Sharing contingency}
\label{sub:we-sharing}

The received contingency, held inside one person, is not yet the beauty relational aesthetics names; the third step carries it into the between. The practice is the timely and unelaborated joint indication of what has been received, the brief look at that light, the few words that say only that the other's remark just now was good, the bringing of the contingent into a focus the two now share and each knows to be shared. Everything turns on timeliness and lightness, because the contingent is by nature passing, and a sharing delayed until it can be explained, or weighted with analysis until it is pinned, has let the coupling event slip while preparing to honour it. The practice is the small, swift, unjudging co-indication that opens the recursive common focus before the moment is gone. It fails in two opposite ways, and the mirror named both: the contingency can be hoarded, kept as a private appreciation never carried across, the connoisseur's quiet possession of a beauty he does not share; or it can be over-spoken, the moment crushed under the weight of saying, the light talked about until it has gone out. Between the hoarding and the crushing lies the brief co-indication that lets the between be born.

This is the principal site when a relation has grown estranged. In estrangement the prior two steps may still occur within each person, each perceiving and even receiving the contingent alone, but the carrying-across has stopped, the field has lost the coupling event, and beauty is produced privately on both sides of a gap it no longer crosses. The labour that matters then is the carrying itself, however small, the re-opening of the channel by which what each sees alone can become a thing the two hold together, because it is only there, in the resumed between, that the beauty estrangement has driven into two private interiors can again become relational.

\subsection{Reproducing contingency}
\label{sub:we-reproducing}

The shared contingency can dissipate, consumed and forgotten, or it can settle so that the next morning's seeing begins higher; the fourth step is the settling. The practice is not the pursuit of a new and stronger contingency but the return to the same daily thing to find in it, this time, a finer appropriateness than before, the same walk at evening, the same greeting, the same dish, met on the hundredth occasion with a perception thickened by the ninety-nine. This is daily renewal made concrete, the freshening of the same rather than the seeking of the new, and it is the practical face of the distinction between the spiral and the wheel: if each return carries a finer registration, the cycle gains and the threshold of joint seeing lowers, so that contingency becomes easier to catch together next time; if each return merely re-enacts the last, the practice has become consumption, and consumption of beauty, like all consumption in the imaginary register, demands ever-larger doses of novelty to feel anything at all. The failure that shadows reproduction is exactly this turn to consumption, the need for a bigger contingency to register a smaller and smaller return, the exhaustion that the series diagnosed as the fate of value sought as a possession.

This is the principal site when a relation has stalled in plenty, when perceiving and receiving and sharing all still happen but produce only a re-enactment of yesterday's pleasure, beauty consumed rather than thickened. And it is here that the worked example closes the loop the model promised, because reproduction's freshening is not the fourth step's solitary achievement but the lowering of a threshold that makes the first step easier on the next turn: the contingency thickened by reproduction is a contingency more readily perceived tomorrow, so that the cycle returns not to where it began but to a perception raised by what the last cycle settled. The capacity this section set out to cultivate, the aesthetic perception of contingency, is therefore not trained by working on perception alone but by turning the whole cycle, since each completed turn lowers the threshold of the perception the next turn begins from. To learn to see what need not have happened is to have received it, shared it, and let it settle, often enough that seeing it has become the easier thing.

\begin{casebox}[The four ways the cultivation of contingency fails]
Each step fails in its own way, and naming the four failures together shows them to be distinct rather than one. Perception fails by the flattened morning: the prior so confident that the unplanned is explained away before it registers, the partner certain she knows how the other is and therefore unable to be surprised by him, the contingency predicted out of existence. Reception fails by the defended morning: the unplanned thing seen but met as interference, corrected back into the plan, the answer that did not fit folded toward the one that would have, the gate raised against what might have revised the hour. Sharing fails in two opposite directions: by hoarding, the contingency kept as a private appreciation never carried across, the quiet connoisseurship of a beauty unshared; and by crushing, the moment over-spoken, weighted with analysis until the light talked about has gone out. Reproduction fails by consumption: the turn to ever-larger novelty, the need for a bigger contingency to feel a smaller return, the beauty had and discarded rather than settled and thickened. Four steps, five failures, no two alike, and the dialectical reading is what tells the practitioner which failure is the one now operating, since the remedy for the flattened morning would do nothing for the defended one, and the remedy for hoarding would only worsen the crushing.
\end{casebox}

\subsection{The shape of the practice}
\label{sub:we-shape}

Set end to end, the four practices are a single small loop that a day affords many times over, and naming it plainly is the point of the section. Before a habitual act, hold the half-second that lets the morning's contingency appear; allow what appears to revise the morning rather than be corrected away; carry it, lightly and at once, into a focus the two of you share; and return, on the next occasion, to find the same thing thickened by the sharing that went before, the threshold of seeing lowered so that the next contingency is caught more readily. The loop closes without returning to its origin, since the last step raises the first, and that closing-without-return is the holonomy the series has tracked from the beginning, here enacted at the smallest scale a shared life contains, the scale of a single noticed thing between two people. None of the four practices requires an hour set aside or a leisure the day does not already hold; each is a different way of meeting an act the day was going to contain regardless, which is the whole of what the next section means when it insists that this cultivation must be available to those whose days hold no spare hours at all. The worked example thus delivers what the theory promised, an aesthetic education that is genuinely practicable, genuinely relational, and genuinely a matter of the everyday, trainable in the most repeated and least spacious hours because it asks not for new hours but for a new way of entering the hours there are.

\subsection{The aptness of the contingency example}
\label{sub:we-why-contingency}

It is worth saying, finally, why contingency was the capacity chosen to walk through the four steps, since the choice was not incidental and the reason illuminates the whole cultivation. Contingency, the noticing of what need not have happened, is the purest case of the perception the paper is about, because it is the perception most directly opposed to the flattening prior that habit installs. The everyday flattens by prediction, by the confident model that explains the morning in advance, and contingency is precisely what the confident model does not predict, the unplanned thing that the flattening prior would absorb without trace. To cultivate the perception of contingency is therefore to cultivate, at its sharpest point, the very capacity to resist the flattening that makes the everyday invisible, and the other perceptions the field-builder needs, the perception of the field's standing state, of the other's mood, of the appropriate gesture, are in a sense easier cases of the same capacity, less sharply opposed to the prior because less purely a matter of the unpredicted. Contingency is the hard case, the perception of exactly what the prior did not contain, and a perception trained on the hard case is trained on the rest, which is why the worked example chose it, as one trains a capacity on its most demanding instance.

Contingency is also the right example because it is the everyday's own source of renewal, the thing that keeps the repeated from becoming the merely repeated. A life that perceived no contingency would be a life wholly predicted, every morning a copy of every morning, the wheel turning without gain; and the perception of contingency is what introduces, into the repeated everyday, the difference that lets the repetition gain rather than recur. The contingency caught and received and shared and let to settle is the difference that makes the hundredth morning more than the first, the small unplanned thing that, woven into the field, keeps the repetition the gaining kind. To cultivate the perception of contingency is thus to cultivate the everyday's capacity to renew itself, the perception of the difference that keeps the repeated from dulling, and this is why the example is not merely one cultivation among many but the cultivation closest to the heart of the whole account, the perception of the contingency that is the everyday's own defense against the wheel. The worked example chose contingency because contingency is where the perception is hardest and where the everyday's renewal is sourced, the sharp point of the capacity and the spring of the good repetition, and a cultivation that masters the perception of contingency has mastered the everyday's resistance to its own flattening, which is as near to the center of the aesthetic education of the everyday as any single capacity comes.
